% MM-15 user manual.  Build with `make manual` (lualatex via latexmk).
% The command reference in Appendix A is generated by `mm_15 --commands-tex`.
\documentclass[11pt,letterpaper]{article}

\usepackage{fontspec}
\IfFontExistsTF{TeX Gyre Pagella}{\setmainfont{TeX Gyre Pagella}}{\setmainfont{Latin Modern Roman}}
\IfFontExistsTF{TeX Gyre Heros}{\setsansfont{TeX Gyre Heros}[Scale=MatchLowercase]}{}
% The stack display uses ℝ ℂ 𝒮. DejaVu Sans Mono has the first two; the script S
% is taken from Latin Modern Math by making the character active inside the
% transcript environments (works with both LuaLaTeX and XeLaTeX).
\IfFontExistsTF{DejaVu Sans Mono}
  {\setmonofont{DejaVu Sans Mono}[Scale=0.86]}
  {\setmonofont{Latin Modern Mono}}
\newfontface\mathglyphs{Latin Modern Math}
\begingroup\catcode`𝒮=\active
\gdef𝒮{{\mathglyphs\char"1D4AE}}\endgroup
\usepackage{unicode-math}
\IfFontExistsTF{TeX Gyre Pagella Math}{\setmathfont{TeX Gyre Pagella Math}}{\setmathfont{Latin Modern Math}}

\usepackage[margin=1.05in]{geometry}
\usepackage{microtype}
\usepackage{xcolor}
\usepackage{booktabs,longtable,array}
\usepackage{fancyvrb}
\usepackage{enumitem}
\usepackage{titlesec}
\usepackage{tcolorbox}
\usepackage{fancyhdr}
\usepackage[hidelinks,pdfusetitle]{hyperref}
\hypersetup{pdftitle={MM-15 User Manual},pdfauthor={Mico Mrkaic}}

\definecolor{codebg}{RGB}{247,247,244}
\definecolor{rule}{RGB}{190,190,185}
\definecolor{accent}{RGB}{120,40,30}

\setlist{nosep}
\titleformat{\section}{\Large\bfseries\sffamily\color{accent}}{\thesection}{0.8em}{}
\titleformat{\subsection}{\large\bfseries\sffamily}{\thesubsection}{0.8em}{}
\titleformat{\subsubsection}{\bfseries\sffamily}{\thesubsubsection}{0.8em}{}
\setcounter{secnumdepth}{2}
\setcounter{tocdepth}{2}

% Session transcripts and scripts.
\DefineVerbatimEnvironment{rpn}{Verbatim}{frame=single,rulecolor=\color{rule},framesep=4pt,fontsize=\small,xleftmargin=0pt,codes={\catcode`𝒮=\active}}
\DefineVerbatimEnvironment{shell}{Verbatim}{frame=leftline,rulecolor=\color{rule},framesep=4pt,fontsize=\small,xleftmargin=0pt}
\newcommand{\code}[1]{\texttt{#1}}
\newcommand{\g}[1]{{\mathglyphs #1}}   % stack type glyphs in running text
\newcommand{\cmd}[1]{\texttt{\textbf{#1}}}
\newcommand{\stk}[1]{\texttt{#1}}
\newenvironment{note}{\begin{tcolorbox}[colback=codebg,colframe=rule,boxrule=0.4pt,arc=1pt,left=6pt,right=6pt,top=4pt,bottom=4pt]}{\end{tcolorbox}}

\pagestyle{fancy}
\fancyhf{}
\fancyhead[L]{\small\sffamily MM-15 User Manual}
\fancyhead[R]{\small\sffamily\nouppercase{\leftmark}}
\fancyfoot[C]{\small\thepage}
\renewcommand{\headrulewidth}{0.3pt}

\title{\sffamily\Huge MM-15\\[0.3em]\Large A programmable RPN calculator for real and complex\\scalars, strings and matrices\\[1em]\large User Manual, version 2.1}
\author{Mico Mrkaic}
\date{November 2026}

\begin{document}
\maketitle
\thispagestyle{empty}

\begin{center}
\begin{minipage}{0.8\textwidth}
\small
MM-15 is free software under the GNU General Public License, version 3 or later.
Source, issues and releases: \url{https://github.com/micomrkaic/MM-15}.
\end{minipage}
\end{center}

\vfill
\tableofcontents
\clearpage

%======================================================================
\section{Introduction}

MM-15 is a stack-based calculator in the tradition of the HP-41C and of FORTH, written in C17 on top of the GNU Scientific Library and GNU Readline. You type numbers and commands; numbers go on a stack and commands take their arguments from it. Beyond the usual scalar functions it handles complex numbers, strings, and real and complex matrices, and it is programmable in three ways: user-defined \emph{words} with conditionals and loops, numerical methods (integration, root finding, minimisation, differentiation) that take any word as the function, and HP-41 style line programs with labels and jumps.

The design has a few fixed points worth knowing before you start.

\begin{itemize}
  \item \textbf{Everything is a token on a line.} Whitespace separates tokens; a line is evaluated left to right. There is no operator precedence because there are no infix operators.
  \item \textbf{Every token is atomic.} If a command fails---wrong type, singular matrix, division by zero---the stack is put back exactly as it was before that token, the error is printed, and the rest of the line is skipped. You never end up with half a result.
  \item \textbf{One promotion rule.} A real function applied outside its real domain returns a complex number rather than \code{nan}: \code{-1 sqrt} is $(0,1)$, \code{2 asin} is $(1.5708, 1.3170)$, and a real matrix with one such entry becomes a complex matrix.
  \item \textbf{Instant startup, no ceremony.} The calculator starts in a few milliseconds, reads your saved words and configuration from \code{\textasciitilde/.config/mm\_15}, and is equally at home in a terminal session and in a shell pipeline.
\end{itemize}

\subsection{Notation used in this manual}

Commands are shown in \cmd{bold monospace}. Their effect on the stack is described by a \emph{stack picture}, written deepest-first with the top of the stack at the right:
\begin{center}\stk{y x -- y+x}\end{center}
means that the command consumes two items---\code{y} below \code{x}---and leaves their sum. A picture that lists the same name on both sides (\stk{A -- A rows cols}) means the item stays on the stack.

Transcripts show what you type, followed by the calculator's output. In the interactive session the whole stack is printed after every line; in the transcripts we usually show only the relevant part. Lines beginning with \code{\textbackslash} are comments and are ignored by the calculator; we use them to annotate results.

\begin{rpn}
2 3 +            \ 5
(1,2) (3,-1) *   \ (5,5)
\end{rpn}

%======================================================================
\section{Getting started}

\subsection{Building and installing}

MM-15 needs a C17 compiler, GNU make, GNU Readline and the GNU Scientific Library (GSL). On Debian or Ubuntu:
\begin{shell}
sudo apt install build-essential libreadline-dev libgsl-dev
\end{shell}
and on macOS with Homebrew, \code{brew install gsl readline}. Then:
\begin{shell}
git clone https://github.com/micomrkaic/MM-15
cd MM-15
make              # builds bin/release/mm_15
make test         # 21 unit cases + golden files, run under AddressSanitizer
make install      # into ~/.local: bin/mm_15, the man page, this manual, and docs
\end{shell}
\code{make install} also seeds \code{\textasciitilde/.config/mm\_15/} with the default configuration and the shipped macros, without overwriting files you already have. \code{PREFIX=/usr/local make install} installs system-wide; \code{make uninstall} removes exactly what was installed. Make sure \code{\textasciitilde/.local/bin} is on your \code{PATH}.

\subsection{A first session}

\begin{rpn}
$ mm_15
MM_RPN>> 2 3 +
[0] ℝ : 5.000000
MM_RPN>> 4 *
[0] ℝ : 20.000000
MM_RPN>> dup sqrt
[0] ℝ : 20.000000
[1] ℝ : 4.472136
MM_RPN>> q
\end{rpn}

The prompt is \code{MM\_RPN>>}. After each line the whole stack is printed, one item per line, deepest first: the index in brackets, a type glyph (\g{ℝ} real, \g{ℂ} complex, \g{𝒮} string, M\g{ℝ} and M\g{ℂ} real and complex matrices), and the value. Position \code{[0]} is the \emph{bottom} of the stack; the highest index is the top, which is where commands take their arguments.

The REPL has GNU Readline editing: arrow keys recall history (kept in \code{.rpn\_history} between sessions), and \textsc{tab} completes command, word and macro names. \cmd{q} exits, saving the configuration; \cmd{undo} restores the stack as it was before the previous line.

\subsection{Getting help inside the calculator}

\begin{description}[style=nextline,leftmargin=1.5em]
  \item[\cmd{help}] a one-screen summary of the input syntax and the command families.
  \item[\cmd{listfcns}] every command, grouped by category, with its stack picture.
  \item[\cmd{"name" usage}] the stack picture, description and an example for one command---or, for a user word, its definition:
\end{description}
\begin{rpn}
MM_RPN>> "mod" usage
"mod" usage: y x -- y mod x   [Arithmetic]
    floored modulo (result has the sign of x)
    Example: -7 3 mod
\end{rpn}
Outside the calculator there is \code{man mm\_15}, and \code{mm\_15 --commands} prints the complete reference as Markdown (it is also Appendix~\ref{app:commands} of this manual).

%======================================================================
\section{Entering data}

All input is case sensitive. A line is split into tokens at whitespace; the exceptions are strings, complex literals and matrix literals, which may contain spaces.

\subsection{Numbers}

Real numbers in the usual forms: \code{3}, \code{-1.25}, \code{.5}, \code{1e3}, \code{-2.5e-3}. A leading \code{+} is accepted. Numbers are IEEE double precision; \cmd{inf} and \cmd{nan} push the special values, and \cmd{pi}, \cmd{e} and \cmd{gravity} ($9.80665\,\mathrm{m/s^2}$) push constants.

\subsection{Complex numbers}

A complex number is written in parentheses, real part first: \code{(1,2)}, \code{(-1.2e-4, 0.7e2)}. Spaces inside the parentheses are allowed. A complex number prints as \code{(1.000000, 2.000000i)}. The commands \cmd{re}, \cmd{im}, \cmd{abs}, \cmd{arg}, \cmd{conj}, \cmd{re2c}, \cmd{split\_c} and \cmd{j2r} convert between the representations:
\begin{rpn}
(3,4) abs        \ 5
(3,4) arg        \ 0.927295  (radians)
(1,2) split_c    \ 1 2       (two reals)
1 2 j2r          \ (1,2)     (imaginary part on top)
5 re2c           \ (5,0)
\end{rpn}
Reals and complex numbers mix freely in arithmetic; the result is complex whenever an operand is.

\subsection{Strings}

Strings are double-quoted: \code{"hello world"}. The escapes \code{\textbackslash"}, \code{\textbackslash\textbackslash}, \code{\textbackslash n} and \code{\textbackslash t} are recognised. Strings are used for dates, file names, word names (\code{"sq" usage}), and as input to \cmd{eval}; see Section~\ref{sec:strings}.

\subsection{Matrices}

An inline matrix gives the dimensions, a dollar sign, and the elements in row-major order:
\begin{rpn}
[2 2 $ 1 2 3 4]             \ the 2x2 matrix  1 2 / 3 4
[1 3 $ 1 2 3]               \ a row vector
[3 1 $ 1 2 3]               \ a column vector
[2 2 $ 1 (0,1) 0 1]         \ one complex entry makes a complex matrix
\end{rpn}
A matrix can also be read from a whitespace-separated text file, giving its shape:
\begin{rpn}
[2, 2, "m.txt"]              \ reads four numbers from m.txt, row by row
\end{rpn}
The stack shows a matrix as \code{M\g{ℝ}: 2 x 2 matrix}; \cmd{pm} (or \cmd{print}) prints it in full. Section~\ref{sec:matrices} covers what you can do with them.

\subsection{Comments}

\code{\textbackslash} and \code{\#} start a comment that runs to the end of the line. This works in the REPL, in scripts, in word definitions and in program files.

\subsection{Execution tokens}

\code{@name} does not run \code{name}; it pushes a reference to it, an \emph{execution token}, as a real number. The numerical methods (Section~\ref{sec:numerics}) and \cmd{exec} consume these: \code{0 pi @sin integrate} integrates the sine over $[0,\pi]$. Any command, user word or macro can be referenced this way.

\subsection{Shell commands}

A line that starts with \code{!} is run by the shell and its standard output is pushed as a string:
\begin{rpn}
!date +%Y
\end{rpn}

%======================================================================
\section{The stack}

The stack holds up to 100 items of any type. The classic FORTH words manipulate it:

\begin{center}
\begin{tabular}{@{}llp{0.55\textwidth}@{}}
\toprule
\cmd{dup}  & \stk{x -- x x}     & duplicate the top \\
\cmd{drop} & \stk{x --}         & discard the top \\
\cmd{swap} & \stk{a b -- b a}   & exchange the top two \\
\cmd{over} & \stk{a b -- a b a} & copy the second item to the top \\
\cmd{nip}  & \stk{a b -- b}     & discard the second item \\
\cmd{tuck} & \stk{a b -- b a b} & copy the top below the second item \\
\cmd{pick} & \stk{... n -- ... v} & copy the item at depth $n$ to the top (\code{0 pick} is \cmd{dup}, \code{1 pick} is \cmd{over}) \\
\cmd{roll} & \stk{... n -- ...}   & move the item at depth $n$ to the top (\code{1 roll} is \cmd{swap}, \code{2 roll} is FORTH's \code{rot}) \\
\cmd{depth} & \stk{-- n}        & number of items on the stack \\
\cmd{clst} & \stk{... --}       & clear the stack \\
\bottomrule
\end{tabular}
\end{center}

\begin{rpn}
10 20 30 2 roll     \ 20 30 10
10 20 30 1 pick     \ 10 20 30 20
\end{rpn}

\cmd{savestack} writes the whole stack to \code{stack\_backup.dat} in the configuration directory and \cmd{loadstack} pushes it back---handy for carrying a set of matrices from one session to the next.

\subsection{Display}

\cmd{print} prints the top item (a matrix in full), \cmd{ps} prints the whole stack, \cmd{pm} prints the matrix on top in full. All three leave the stack unchanged.

Numbers are shown in fixed-point with six decimals by default. \cmd{setprec} sets the number of decimals (0--20), and \cmd{sfs} toggles between fixed-point and scientific notation, where the precision is the number of significant digits:
\begin{rpn}
10 setprec 1 3 /      \ 0.3333333333
sfs 1 3 /             \ 0.3333333333  (6 significant digits would be 0.333333)
3 setprec 12345.678   \ 1.23e+04
sfs                   \ back to fixed point
\end{rpn}
Both settings are saved in \code{config.txt} when you quit.

%======================================================================
\section{Arithmetic and functions}

\subsection{Basic arithmetic}

\cmd{+}, \cmd{-}, \cmd{*}, \cmd{/} and \cmd{\textasciicircum} (also spelled \cmd{pow}) work on reals, complex numbers and matrices (Section~\ref{sec:matrices} explains the matrix cases). Division by zero is an error, not \code{inf}.

\begin{center}
\begin{tabular}{@{}llp{0.55\textwidth}@{}}
\toprule
\cmd{chs}  & \stk{x -- -x}   & change sign \\
\cmd{inv}  & \stk{x -- 1/x}  & reciprocal \\
\cmd{abs}  & \stk{z -- |z|}  & absolute value or complex magnitude \\
\cmd{sqrt} & \stk{x -- sqrt(x)} & square root \\
\cmd{intg} & \stk{x -- int(x)}  & integer part, truncated toward zero (\code{-3.7 intg} is $-3$) \\
\cmd{frac} & \stk{x -- frac(x)} & fractional part, $x - \mathrm{intg}(x)$ (\code{-3.7 frac} is $-0.7$) \\
\cmd{mod}  & \stk{y x -- y mod x} & floored modulo; the result has the sign of $x$ (\code{-7 3 mod} is 2) \\
\cmd{fact} & \stk{x -- x!}   & factorial, $\Gamma(x+1)$ for non-integers \\
\cmd{ncr}  & \stk{n k -- C(n,k)} & combinations \\
\cmd{npr}  & \stk{n k -- P(n,k)} & permutations \\
\bottomrule
\end{tabular}
\end{center}

The three percentage functions follow the HP convention of \emph{keeping} the base value $y$ so that it is available for the next step:
\begin{rpn}
200 25 %        \ 200 50      25% of 200
200 50 %t       \ 200 25      50 is 25% of 200
80 100 %chg     \ 80 25       from 80 to 100 is +25%
\end{rpn}

\subsection{Exponentials, logarithms, trigonometry}

\cmd{exp}, \cmd{ln}, \cmd{log} (base 10); \cmd{sin}, \cmd{cos}, \cmd{tan}, \cmd{asin}, \cmd{acos}, \cmd{atan}; \cmd{sinh}, \cmd{cosh}, \cmd{tanh}, \cmd{asinh}, \cmd{acosh}, \cmd{atanh}. Angles are in radians. All of them accept real and complex arguments and apply elementwise to matrices.

\subsection{The promotion rule}
\label{sec:promotion}

Real functions whose result would not be real return a complex number. The rule is applied consistently, so a calculation never produces \code{nan} where a complex answer exists:

\begin{center}
\begin{tabular}{@{}lll@{}}
\toprule
Command & Promotes when & Example \\
\midrule
\cmd{sqrt}  & $x < 0$   & \code{-4 sqrt} $\to (0, 2)$ \\
\cmd{ln}, \cmd{log} & $x < 0$ & \code{-1 ln} $\to (0, \pi)$ \\
\cmd{asin}, \cmd{acos} & $|x| > 1$ & \code{2 asin} $\to (1.5708, 1.3170)$ \\
\cmd{acosh} & $x < 1$   & \code{0.5 acosh} $\to (0, 1.0472)$ \\
\cmd{atanh} & $|x| > 1$ & \code{2 atanh} $\to (0.5493, 1.5708)$ \\
\cmd{\textasciicircum} & $y < 0$ and $x$ not an integer & \code{-8 3 inv \textasciicircum} $\to (1, 1.7321)$ \\
\bottomrule
\end{tabular}
\end{center}

A real matrix with one such entry becomes a complex matrix: \code{[2 2 \$ 1 -2 3 4] ln} is complex. The limits $\pm\infty$ remain real where they are the correct answer (\code{0 ln} is \code{-inf}).

\subsection{Special functions and the normal distribution}

\cmd{gamma}, \cmd{ln\_gamma}, \cmd{beta} and \cmd{ln\_beta} wrap the GSL special functions. \cmd{npdf}, \cmd{ncdf} and \cmd{nquant} are the standard normal density, distribution function and quantile:
\begin{rpn}
1.96 ncdf        \ 0.975002
0.975 nquant     \ 1.959964
5 gamma          \ 24
2 3 beta         \ 0.083333
\end{rpn}

\subsection{Comparisons and logic}

\cmd{eq}, \cmd{neq}, \cmd{lt}, \cmd{leq}, \cmd{gt}, \cmd{geq} leave a flag: 1 for true, 0 for false. \cmd{and}, \cmd{or} and \cmd{not} treat 0 as false and anything else as true. Complex numbers compare by magnitude; strings support \cmd{eq} and \cmd{neq}; matrices compare elementwise and produce a matrix of flags, and a scalar compared with a matrix is compared with every element.
\begin{rpn}
2 3 lt                     \ 1
"abc" "abc" eq             \ 1
[1 3 $ 1 5 3] 2 gt pm      \ | 0 1 1 |
\end{rpn}
Flags are what the control words inside user words test (Section~\ref{sec:words}).

%======================================================================
\section{Matrices}
\label{sec:matrices}

\subsection{Creating matrices}

Besides the literal forms of Section~3, there are constructors:
\begin{rpn}
3 eye                      \ 3x3 identity
2 3 ones                   \ 2x3 matrix of ones
2 3 zeroes                 \ 2x3 matrix of zeros
2 2 rand                   \ uniform(0,1) entries
2 2 randn                  \ standard normal entries
0 0.25 1 rrange            \ row vector 0 0.25 0.5 0.75 1
[1 3 $ 1 2 3] to_diag      \ 3x3 diagonal matrix with 1 2 3 on the diagonal
\end{rpn}

\subsection{Arithmetic}

\begin{itemize}
  \item \cmd{+} and \cmd{-} add and subtract matrices of the same shape. A scalar is added to every element.
  \item \cmd{*} is the matrix product; a scalar scales every element.
  \item \cmd{/} divides a matrix by a scalar elementwise; \code{x A /} is $x/a_{ij}$ elementwise; and \code{A B /} is $A B^{-1}$ (right division).
  \item \cmd{\textasciicircum} raises a square matrix to a non-negative integer power ($A^0$ is the identity).
  \item \cmd{.*}, \cmd{./}, \cmd{.\textasciicircum} are the elementwise product, quotient and power; either operand may be a scalar.
  \item A real matrix combined with a complex matrix or complex scalar is promoted to complex first.
\end{itemize}
\begin{rpn}
[2 2 $ 1 2 3 4] dup * pm       \ |  7 10 |  matrix product
                               \ | 15 22 |
[2 2 $ 1 2 3 4] dup .* pm      \ |  1  4 |  elementwise
                               \ |  9 16 |
[2 2 $ 1 2 3 4] 10 + pm        \ | 11 12 | / | 13 14 |
[2 2 $ 1 2 3 4] 3 ^ pm         \ | 37 54 | / | 81 118 |
\end{rpn}

The unary functions---\cmd{sqrt}, \cmd{exp}, \cmd{ln}, the trigonometric and hyperbolic functions, \cmd{abs}, \cmd{chs}, \cmd{inv}, \cmd{intg}, \cmd{frac}, \cmd{npdf} and the rest---apply elementwise. Note that \cmd{inv} is therefore the elementwise reciprocal; the matrix inverse is \cmd{minv}.

\subsection{Linear algebra}

\begin{center}
\begin{tabular}{@{}llp{0.6\textwidth}@{}}
\toprule
\cmd{tran}, \cmd{'} & \stk{A -- A\textasciicircum T} & transpose \\
\cmd{minv} & \stk{A -- A\textasciicircum -1} & inverse of a square real or complex matrix; singular is an error \\
\cmd{pinv} & \stk{A -- A\textasciicircum +}  & Moore--Penrose pseudoinverse (real, square), by SVD \\
\cmd{det}  & \stk{A -- det(A)} & determinant \\
\cmd{eig}  & \stk{A -- V L}    & eigenvectors (columns of $V$) and eigenvalues (diagonal of $L$), both complex \\
\cmd{chol} & \stk{A -- L}      & Cholesky factor of a symmetric positive definite matrix, $A = LL^T$ \\
\cmd{svd}  & \stk{A -- U S V}  & singular value decomposition, $A = U S V^T$ \\
\cmd{kron} & \stk{A B -- A(x)B} & Kronecker product \\
\bottomrule
\end{tabular}
\end{center}
\begin{rpn}
[2 2 $ 4 7 2 6] det            \ 10
[2 2 $ 4 7 2 6] minv pm        \ |  0.6 -0.7 |
                               \ | -0.2  0.4 |
[2 2 $ 2 1 1 2] eig pm         \ L: eigenvalues 3 and 1 on the diagonal
[2 2 $ 4 2 2 3] chol pm        \ | 2 0 | / | 1 1.4142 |
\end{rpn}
Solving $Ax = b$ is \code{b A minv *}---or, keeping the operands in the natural order, \code{A b swap minv *}. The shipped macros add \code{cond} (condition number), \code{norm2}, \code{normfro} and \code{XprimeX} (Appendix~\ref{app:macros}).

\subsection{Elements and shape}

Indices are 0-based.
\begin{center}
\begin{tabular}{@{}llp{0.55\textwidth}@{}}
\toprule
\cmd{dim}       & \stk{A -- A rows cols} & dimensions; the matrix stays \\
\cmd{get\_aij}  & \stk{A i j -- A a\_ij}  & element $(i,j)$; the matrix stays \\
\cmd{set\_aij}  & \stk{A i j x -- A}     & set element $(i,j)$ in place; a complex $x$ promotes a real matrix \\
\cmd{reshape}   & \stk{A rows cols -- B} & reshape in row-major order \\
\cmd{split\_mat} & \stk{A -- a11 a12 ...} & push every element as a scalar, row by row \\
\cmd{diag}      & \stk{A -- v}    & the main diagonal as a row vector \\
\cmd{to\_diag}  & \stk{v -- D}    & diagonal matrix from a row or column vector \\
\cmd{join\_v}   & \stk{A B -- [A; B]} & stack vertically (same number of columns) \\
\cmd{join\_h}   & \stk{A B -- [A B]}  & join side by side (same number of rows) \\
\bottomrule
\end{tabular}
\end{center}
\begin{rpn}
[2 2 $ 1 2 3 4] 0 1 get_aij        \ A 2
[2 2 $ 1 2 3 4] 1 0 9 set_aij pm   \ | 1 2 | / | 9 4 |
[2 3 $ 1 2 3 4 5 6] 3 2 reshape pm \ | 1 2 | / | 3 4 | / | 5 6 |
[1 2 $ 1 2] [1 2 $ 3 4] join_v pm  \ | 1 2 | / | 3 4 |
\end{rpn}

\subsection{Reductions and statistics}

The \code{c}-prefixed commands reduce each column and return a row vector; the \code{r}-prefixed ones reduce each row and return a column vector. They consume the matrix: \cmd{csum}/\cmd{rsum}, \cmd{cmean}/\cmd{rmean}, \cmd{cvar}/\cmd{rvar} (sample variance, $n-1$ in the denominator), \cmd{cmin}/\cmd{rmin}, \cmd{cmax}/\cmd{rmax}. \cmd{cumsum\_r} and \cmd{cumsum\_c} are cumulative sums along rows and down columns; the macros \code{cumprod\_r} and \code{cumprod\_c} do the same for products.
\begin{rpn}
[2 2 $ 1 2 3 4] cmean pm       \ | 2 3 |
[2 2 $ 1 2 3 4] rsum pm        \ | 3 | / | 7 |
[2 2 $ 1 2 3 4] cvar pm        \ | 2 2 |
\end{rpn}

\subsection{Polynomials}

A polynomial is a vector of coefficients, \emph{highest degree first}, as in Matlab. \cmd{pval} evaluates it at a real or complex point; \cmd{roots} returns all roots as a complex row vector.
\begin{rpn}
[1 3 $ 1 -3 2] 2 pval          \ 0        x^2 - 3x + 2 at x = 2
[1 3 $ 1 -3 2] roots pm        \ | (1,0) (2,0) |
[1 3 $ 1 0 1] roots pm         \ | (0,1) (0,-1) |   x^2 + 1
\end{rpn}
The macros \code{npv} and \code{irr} are built on these two commands (Appendix~\ref{app:macros}).

%======================================================================
\section{Strings}
\label{sec:strings}

\begin{center}
\begin{tabular}{@{}llp{0.55\textwidth}@{}}
\toprule
\cmd{scon}    & \stk{s1 s2 -- s1s2} & concatenate \\
\cmd{substr}  & \stk{s start end -- sub} & substring, 0-based, end exclusive \\
\cmd{slen}    & \stk{s -- n}  & length in bytes \\
\cmd{s2u}, \cmd{s2l} & \stk{s -- s} & upper and lower case \\
\cmd{srev}    & \stk{s -- s}  & reverse \\
\cmd{int2str} & \stk{x -- s}  & integer part of $x$ as a string \\
\cmd{eval}    & \stk{s --}    & evaluate the string as a line of input \\
\bottomrule
\end{tabular}
\end{center}
\begin{rpn}
"abcdef" 1 3 substr        \ "bc"
"foo" "bar" scon s2u       \ "FOOBAR"
"2 3 +" eval               \ 5
"12:45" dec_hr             \ 12.75   (a shipped macro: substr + eval)
\end{rpn}
\cmd{eval} is the bridge from text to computation: the macro \code{dec\_hr} converts a time string such as the one \cmd{sunrise} returns into decimal hours by cutting out the hour and minute fields with \cmd{substr} and evaluating them.

%======================================================================
\section{Registers}

There are 64 storage registers, numbered 0--63, and each can hold any type---a matrix or a string as well as a number. Following the HP convention the value comes first and the register number last:
\begin{rpn}
42 0 sto             \ store 42 in R0
0 rcl                \ 42   (a copy; the register keeps its value)
[2 2 $ 1 2 3 4] 5 sto
pr                   \ shows which registers are occupied
ffr                  \ 1    index of the first free register
\end{rpn}
The macros \code{sto\_add}, \code{sto\_sub}, \code{sto\_mul} and \code{sto\_div} are the HP storage arithmetic: \code{5 0 sto\_add} adds 5 to register 0. \cmd{saveregs} and \cmd{loadregs} write and read \code{registers.txt} in the configuration directory; \cmd{clregs} empties all registers.

Registers are the natural way to pass parameters to a word used in a numerical method (Section~\ref{sec:numerics}): the word reads the parameter with \cmd{rcl} and you change it with \cmd{sto} between calls.

%======================================================================
\section{Words: programming the calculator}
\label{sec:words}

A \emph{word} is a named sequence of tokens. Defining one takes a single line that starts with a colon and ends with a semicolon:
\begin{rpn}
: sq dup * ;
5 sq                 \ 25
: hyp sq swap sq + sqrt ;
3 4 hyp              \ 5
\end{rpn}
Once defined, a word is used exactly like a built-in command; it can appear in other words, in scripts, in programs, and as a function for the numerical methods. \cmd{listwords} lists your words with their definitions, \code{"hyp" usage} shows one, \code{"hyp" delword} deletes one, and \cmd{clrwords} deletes them all.

Words are \emph{compiled} when defined. Every name in the body must already exist (as a command, a word or a macro), so a typo is reported immediately rather than when the word is first run; the control structures described below are turned into branches at this point. A word may call itself, and nesting stops with an error at 256 levels. Redefining a word replaces it, and any word that calls it sees the new definition.

A name must not be a built-in command or control word and must not start with a digit. Names are case sensitive.

\subsection{Conditionals}

Inside a word, \cmd{if} consumes a flag from the stack and executes the following tokens when it is true (non-zero); \cmd{else} introduces the alternative; \cmd{then} closes the construct.
\begin{rpn}
: abs2 dup 0 lt if chs then ;
-3 abs2              \ 3
: sign dup 0 lt if drop -1 else 0 gt if 1 else 0 then then ;
-7 sign 0 sign 2 sign     \ -1 0 1
\end{rpn}

\subsection{Loops}

\begin{description}[style=nextline,leftmargin=1.5em]
  \item[\code{begin ... flag until}] repeat until the flag is true (the body runs at least once).
  \item[\code{begin ... flag while ... repeat}] test in the middle; leave the loop when the flag is false.
  \item[\code{limit start do ... loop}] a counted loop from \code{start} up to but not including \code{limit}; \cmd{i} is the current index and \cmd{j} the index of the enclosing loop. If \code{start} equals \code{limit} the body is skipped.
  \item[\code{limit start do ... n +loop}] the same with a step of $n$; a negative step counts down to \code{limit} inclusive.
\end{description}
Loops nest to 16 levels. Flags are reals; 0 is false.
\begin{rpn}
: fac dup 1 gt if dup 1 - fac * then ;         \ recursive
10 fac                                       \ 3628800

: sum 0 swap 1 + 1 do i + loop ;             \ 1 + 2 + ... + n
100 sum                                      \ 5050

: step dup 2 mod if 3 * 1 + else 2 / then ;   \ one Collatz step
: collatz 0 swap begin dup 1 gt while swap 1 + swap step repeat drop ;
27 collatz                                   \ 111 steps

: fib 0 swap 1 swap 0 do tuck + loop drop ;  \ nth Fibonacci number
30 fib                                       \ 832040

: table 1 + 1 do i dup sq loop ;             \ i and i^2 for i = 1..n
3 table                                      \ 1 1 2 4 3 9
\end{rpn}
A definition occupies the whole line: nothing may follow the closing semicolon. Long bodies are best split into words, as \code{collatz} does with \code{step}.

\subsection{Saving words}

\cmd{savewords} writes all your words to \code{user\_words.txt} in the configuration directory, one \code{name body} per line, and they are loaded automatically at the next start; \cmd{loadwords} reloads the file by hand. The file is plain text and you may edit it: forward references are fine (it is loaded in two passes), comments start with \code{\textbackslash}, and a line that fails to compile is reported with its line number and skipped.

\subsection{Macros}

Macros are words that ship with the calculator in \code{predefined\_macros.txt}, loaded at startup from the configuration directory; \cmd{listmacros} lists them. They behave exactly like words---the only difference is where they are stored, so that \cmd{clrwords} does not remove them. Appendix~\ref{app:macros} describes each. If a word and a macro have the same name, the word wins. \code{make install} replaces the macro file with the version that matches the program, so put your own definitions in \code{user\_words.txt}, not there.

%======================================================================
\section{Numerical methods}
\label{sec:numerics}

The numerical methods take a function as an execution token---\code{@name} for any word or command---and treat it as a map $x \mapsto f(x)$: the function is called with $x$ on the stack and must leave exactly one real number. An error inside the function aborts the method, and the stack is restored.

\begin{center}
\begin{tabular}{@{}llp{0.55\textwidth}@{}}
\toprule
\cmd{integrate} & \stk{a b xt -- I} & $\int_a^b f$ by tanh--sinh quadrature with adaptive bisection \\
\cmd{romberg}   & \stk{a b xt -- I} & $\int_a^b f$ by Romberg extrapolation \\
\cmd{fzero}     & \stk{a b xt -- x} & root in a bracket $[a,b]$ where $f$ changes sign (Brent) \\
\cmd{solve}     & \stk{g xt -- x}   & root near one guess $g$ (HP-15C style: expand a bracket, then Brent) \\
\cmd{minimize}  & \stk{a b xt -- x} & argument of the minimum of $f$ on $[a,b]$ (Brent) \\
\cmd{deriv}     & \stk{x xt -- f'(x)} & derivative by Ridders' extrapolation \\
\cmd{fsum}, \cmd{fprod} & \stk{lo hi xt -- S} & $\sum f(i)$ or $\prod f(i)$ over the integers $lo..hi$ \\
\cmd{exec}      & \stk{xt --}       & run the token \\
\cmd{tol}       & \stk{x --}        & set the tolerance, relative to $\max(1,|\text{result}|)$; default $10^{-10}$ \\
\cmd{nevals}    & \stk{-- n}        & function evaluations used by the last call \\
\cmd{interr}    & \stk{-- e}        & error estimate of the last \cmd{integrate} \\
\bottomrule
\end{tabular}
\end{center}

\begin{rpn}
12 setprec
0 2 @exp integrate              \ 6.389056098931   (e^2 - 1)
: rsqrt sqrt inv ;
0 1 @rsqrt integrate            \ 2.000000000000   endpoint singularity, exact
0 pi @sin integrate             \ 2.000000000000
0 1 @exp romberg                \ 1.718281828459
2 4 @sin fzero                  \ 3.141592653605
3 @sin solve                    \ 3.141592653590
: par 3 - dup * 1 + ;
0 5 @par minimize               \ 3.000000000000
2 @exp deriv                    \ 7.389056098931
: invsq dup * inv ;
1 1000 @invsq fsum              \ 1.643934566682   (pi^2/6 - 1/1000 + ...)
3 @par exec                     \ 1
\end{rpn}

Some practical notes.
\begin{itemize}
  \item \cmd{integrate} uses the tanh--sinh rule, whose nodes cluster at the endpoints; an integrable singularity \emph{at} an endpoint converges to full precision, and the method is exponentially convergent for analytic integrands. Interior kinks or peaks are handled by bisecting the interval where the rule does not converge. If the error estimate stays above the tolerance a warning is printed and the result is returned anyway; \cmd{interr} tells you how good it is.
  \item \cmd{fzero} needs a sign change; \cmd{solve} needs only a guess and will grow a bracket around it, so it is the one to reach for interactively. Both are exact to the tolerance for simple roots.
  \item A function may take parameters from registers: \code{: f 0 rcl * sin ;} integrates $\sin(kx)$ with $k$ in register 0, and \code{2 0 sto} changes $k$.
  \item The function may be a built-in (\code{@sin}, \code{@exp}) or a macro as well as a word. It may also be a long-running word; the limit is two million evaluations per call, after which the method gives up with an error suggesting a looser tolerance.
  \item Loosen \cmd{tol} for expensive functions: \code{1e-6 tol}. The value is kept in \code{config.txt}.
\end{itemize}
The methods were cross-checked against SciPy on 62 cases with a worst relative error of $3.5\times10^{-11}$.

%======================================================================
\section{Dates, time and the sun}

Dates are strings in the form \code{d.m.yyyy}; leading zeros are optional on input and present on output (\code{"1.2.2025"} and \code{"01.02.2025"} are the same date). Times are strings \code{HH:MM}.

\begin{center}
\begin{tabular}{@{}llp{0.5\textwidth}@{}}
\toprule
\cmd{today}     & \stk{-- date}      & today's date \\
\cmd{now}       & \stk{-- time}      & current local time \\
\cmd{dateplus}  & \stk{date n -- date} & $n$ days later ($n$ may be negative) \\
\cmd{ddays}     & \stk{d1 d2 -- n}   & days from \code{d1} to \code{d2} \\
\cmd{dow}       & \stk{date -- name} & weekday name \\
\cmd{edmy}      & \stk{date -- y m d} & split into year, month, day \\
\cmd{num2date}  & \stk{d m y -- date} & build a date \\
\cmd{days2eoy}  & \stk{date -- n}    & days from the date to the end of its year \\
\cmd{tz\_offset} & \stk{date zone -- h} & UTC offset in hours of an IANA zone on that date \\
\cmd{sunrise}, \cmd{sunset} & \stk{date lat lon zone -- time} & local sunrise and sunset \\
\cmd{dawn}, \cmd{dusk} & \stk{date lat lon zone -- time} & civil dawn and dusk (sun $6^\circ$ below the horizon) \\
\bottomrule
\end{tabular}
\end{center}
\begin{rpn}
"1.1.2025" 30 dateplus         \ "31.01.2025"
"1.1.2025" "31.12.2025" ddays  \ 364
"1.1.2025" dow                 \ "Wednesday"
"1.3.2024" days2eoy            \ 305
"1.12.2026" "America/New_York" tz_offset        \ -5
"21.6.2026" 46.24 14.36 "Europe/Ljubljana" sunset   \ "20:58"
\end{rpn}
The \code{zone} argument of the solar functions is either an IANA zone name such as \code{"Europe/Ljubljana"} or a UTC offset in hours such as \code{-4}. With a zone name the offset is resolved \emph{for the date on the stack}, so daylight-saving time is never your problem; names are validated against the system's \code{zoneinfo} database. The shipped location macros push latitude, longitude and zone in one go:
\begin{rpn}
today 100 dateplus durham dawn      \ civil dawn in Durham, NC, 100 days from now
today kranj sunset dec_hr           \ sunset in Kranj as a decimal hour
\end{rpn}
Add your own with one line in \code{user\_words.txt}: \code{home 51.5 -0.12 "Europe/London"}.

%======================================================================
\section{Programs}

Besides words, MM-15 runs HP-41 style \emph{programs}: text files with one instruction per line, labels, jumps and subroutine calls. They are useful for longer procedures with explicit control flow, and they are how the calculator was programmed before words had loops.

\begin{description}[style=nextline,leftmargin=1.5em]
  \item[\code{LBL name}] a label.
  \item[\code{GOTO name}, \code{GOSUB name}, \code{RTN}] jump, call, and return.
  \item[\code{END}] stop (the program also stops at the end of the file).
  \item[\emph{tests}] a line consisting of a test such as \code{top\_gt?} skips the next line when the test is false. The tests do not consume anything. With $y$ below $x$ on the stack: \code{top\_eq?}, \code{top\_neq?}, \code{top\_lt?}, \code{top\_leq?}, \code{top\_gt?}, \code{top\_gte?} compare $y$ with $x$; \code{top\_eq0?}, \code{top\_neq0?}, \code{top\_lt0?}, \code{top\_lte0?}, \code{top\_gt0?}, \code{top\_gte0?} compare $x$ with zero.
  \item[anything else] is evaluated as a line of ordinary input, so a line may hold several tokens.
\end{description}
Blank lines and lines starting with \code{\textbackslash} or \code{\#} are ignored.

A factorial, \code{fact.prg}:
\begin{rpn}
\ n -- n!
1
LBL loop
over 1
top_gt?
GOTO body
GOTO done
LBL body
drop *
swap 1 - swap
GOTO loop
LBL done
drop drop nip
END
\end{rpn}
\begin{rpn}
5 "fact.prg" run        \ 120
\end{rpn}
\cmd{run} loads and runs a program file; with \code{verbose\_mode = 1} in the configuration it also prints the listing. Two related commands: \cmd{batch} runs a file of ordinary lines (no labels or tests), and \cmd{eval} evaluates a string.

%======================================================================
\section{Batch mode and scripting}

When MM-15 is given expressions or files on the command line, or when its standard input is not a terminal, it runs in \emph{batch mode}: no banner, no per-line echo, the final stack is printed, and the exit status is 1 if any error occurred.
\begin{shell}
mm_15 -e "2 10 ^"                         # one line
mm_15 -e "[2 2 $ 4 7 2 6] det" -e "sqrt"  # several, in order
mm_15 -f script.rpn                       # a file; '-' means stdin
echo "1 3 /" | mm_15 -n                   # a pipeline
mm_15 -n -e ': f dup * 2 - ;' -e '0 3 @f fzero'
\end{shell}
\code{-n} starts with a clean slate: no saved words, macros or configuration, and nothing is written at exit. Use it in scripts so that their behaviour does not depend on the state of your configuration directory. Scripts may define words, and \code{\textbackslash} or \code{\#} comments are allowed. The project's own tests are batch runs compared against golden output files (\code{tests/golden/}).

\begin{center}
\begin{tabular}{@{}lp{0.7\textwidth}@{}}
\toprule
\code{-e EXPR} & evaluate one line; may be repeated \\
\code{-f FILE} & run the lines of a file; \code{-} is standard input; may be repeated and mixed with \code{-e} \\
\code{-n}      & do not load the saved macros, words and configuration \\
\code{--commands} & print the command reference as Markdown \\
\code{--commands-tex} & the same as LaTeX tables (used to build this manual) \\
\code{--version}, \code{-h} & version; usage \\
\bottomrule
\end{tabular}
\end{center}
Exit status: 0 on success, 1 if any error was reported in batch mode, 2 for a usage error.

%======================================================================
\section{Configuration and files}

Everything lives in \code{\$XDG\_CONFIG\_HOME/mm\_15/}, which defaults to \code{\textasciitilde/.config/mm\_15/}:

\begin{center}
\begin{tabular}{@{}lp{0.65\textwidth}@{}}
\toprule
\code{config.txt} & \code{print\_precision}, \code{fixed\_point}, \code{verbose\_mode}, \code{tol}; written at exit \\
\code{predefined\_macros.txt} & macros, one \code{name body} per line; replaced by \code{make install} \\
\code{user\_words.txt} & words saved by \cmd{savewords}; loaded at startup \\
\code{registers.txt} & \cmd{saveregs} / \cmd{loadregs} \\
\code{stack\_backup.dat} & \cmd{savestack} / \cmd{loadstack} \\
\code{.rpn\_history} & readline history (last 1000 lines) \\
\bottomrule
\end{tabular}
\end{center}
A missing configuration directory is created on first use. \code{verbose\_mode = 1} makes the calculator chattier: it lists macros at startup, echoes word definitions in batch mode, and prints program listings.

The configuration directory is per user; the \code{data/} directory in the source tree holds the defaults that \code{make install} copies there if nothing exists yet.

%======================================================================
\section{Errors}

Every top-level token is atomic. When a command fails, the stack is restored to its state before that token, a message beginning with \code{Error:} is printed on standard error, and the rest of the line is skipped:
\begin{rpn}
MM_RPN>> 1 2 [2 2 $ 1 2 2 4] minv +
Error: Matrix is singular, cannot invert
[0] ℝ : 1.000000
[1] ℝ : 2.000000
[2] Mℝ: 2 x 2 matrix
\end{rpn}
The \cmd{+} after the failing \cmd{minv} did not run. A user word counts as one token: if it fails half-way, nothing it did remains. \cmd{undo} goes one step further back and restores the stack as it was before the previous line.

Common messages:
\begin{description}[leftmargin=1.5em]
  \item[\code{needs N item(s) on the stack}] a command was given too few arguments.
  \item[\code{unsupported operand types}] e.g.\ adding a string to a number, or multiplying matrices of incompatible shapes.
  \item[\code{Division by zero}] from \cmd{/}, \cmd{inv} or \cmd{./}.
  \item[\code{must leave exactly one real result}] a function passed to a numerical method pushed nothing, several items, or a non-real.
  \item[\code{runaway recursion}] a word called itself more than 256 levels deep.
\end{description}
In batch mode any error makes the exit status 1, so a script can be checked with \code{\&\&} like any other command.

%======================================================================
\appendix
\clearpage
\section{Command reference}
\label{app:commands}

Generated from the dispatch table by \code{mm\_15 --commands-tex}; the same information is available at run time from \cmd{listfcns} and \code{"name" usage}. Stack pictures read deepest-first.

\small
% Generated by mm_15 --commands-tex from the dispatch table; do not edit by hand.
\subsection*{Arithmetic}
\begin{longtable}{@{}>{\raggedright\arraybackslash}p{0.12\textwidth}>{\raggedright\arraybackslash}p{0.19\textwidth}>{\raggedright\arraybackslash}p{0.36\textwidth}>{\raggedright\arraybackslash}p{0.23\textwidth}@{}}
\toprule \textbf{Command} & \textbf{Stack} & \textbf{Description} & \textbf{Example}\\ \midrule \endfirsthead
\toprule \textbf{Command} & \textbf{Stack} & \textbf{Description} & \textbf{Example}\\ \midrule \endhead
\bottomrule \endlastfoot
\texttt{+} & \texttt{y x -{}-{} y+x} & add scalars or matrices; a scalar added to a matrix is added elementwise & \texttt{3 4 +}\\
\texttt{-{}} & \texttt{y x -{}-{} y-{}x} & subtract scalars or matrices (a scalar elementwise) & \texttt{10 4 -{}}\\
\texttt{*} & \texttt{y x -{}-{} y*x} & multiply scalars, or matrix product for matrices & \texttt{[2 2 \$ 1 2 3 4] [2 1 \$ 1 1] *}\\
\texttt{/} & \texttt{y x -{}-{} y/x} & divide scalars, or a matrix by a scalar; division by zero is an error & \texttt{1 3 /}\\
\texttt{\^{}} & \texttt{y x -{}-{} y\^{}x} & power; a square matrix to an integer power is a matrix power; a negative base with a fractional exponent promotes to complex & \texttt{2 10 \^{}}\\
\texttt{pow} & \texttt{y x -{}-{} y\^{}x} & power (alias of \^{}) & \texttt{2 10 pow}\\
\texttt{.*} & \texttt{A B -{}-{} A.*B} & elementwise matrix product & \texttt{[2 2 \$ 1 2 3 4] [2 2 \$ 1 1 2 2] .*}\\
\texttt{./} & \texttt{A B -{}-{} A./B} & elementwise division (matrix or scalar operands) & \texttt{[2 2 \$ 1 2 3 4] [2 2 \$ 1 2 1 2] ./}\\
\texttt{.\^{}} & \texttt{A B -{}-{} A.\^{}B} & elementwise power (matrix or scalar operands) & \texttt{[2 2 \$ 1 2 3 4] 2 .\^{}}\\
\texttt{chs} & \texttt{x -{}-{} -{}x} & change sign (elementwise for matrices) & \texttt{5 chs}\\
\texttt{inv} & \texttt{x -{}-{} 1/x} & reciprocal (elementwise for matrices; see minv for the matrix inverse) & \texttt{4 inv}\\
\texttt{abs} & \texttt{z -{}-{} \textbar{}z\textbar{}} & absolute value; magnitude for complex; elementwise for matrices & \texttt{(3,4) abs}\\
\texttt{sqrt} & \texttt{x -{}-{} sqrt(x)} & square root; a negative real promotes to complex & \texttt{-{}4 sqrt}\\
\texttt{\%} & \texttt{y x -{}-{} y y*x/100} & percent of y (HP \%; keeps y) & \texttt{200 25 \%}\\
\texttt{\%t} & \texttt{y x -{}-{} y 100*x/y} & x as a percent of the total y (HP \%T; keeps y) & \texttt{200 50 \%t}\\
\texttt{\%chg} & \texttt{y x -{}-{} y 100*(x-{}y)/y} & percent change from y to x (HP delta\%; keeps y) & \texttt{80 100 \%chg}\\
\texttt{mod} & \texttt{y x -{}-{} y mod x} & floored modulo (result has the sign of x) & \texttt{-{}7 3 mod}\\
\texttt{fact} & \texttt{x -{}-{} x!} & factorial; Gamma(x+1) for non-{}integers & \texttt{5 fact}\\
\texttt{ncr} & \texttt{n k -{}-{} C(n,k)} & combinations & \texttt{52 5 ncr}\\
\texttt{npr} & \texttt{n k -{}-{} P(n,k)} & permutations & \texttt{10 3 npr}\\
\texttt{frac} & \texttt{x -{}-{} frac(x)} & fractional part, x -{} intg(x), with the sign of x & \texttt{3.7 frac}\\
\texttt{intg} & \texttt{x -{}-{} int(x)} & integer part, truncated toward zero as on HP calculators (-{}3.7 gives -{}3) & \texttt{3.7 intg}\\
\end{longtable}

\subsection*{Exp/Log}
\begin{longtable}{@{}>{\raggedright\arraybackslash}p{0.12\textwidth}>{\raggedright\arraybackslash}p{0.19\textwidth}>{\raggedright\arraybackslash}p{0.36\textwidth}>{\raggedright\arraybackslash}p{0.23\textwidth}@{}}
\toprule \textbf{Command} & \textbf{Stack} & \textbf{Description} & \textbf{Example}\\ \midrule \endfirsthead
\toprule \textbf{Command} & \textbf{Stack} & \textbf{Description} & \textbf{Example}\\ \midrule \endhead
\bottomrule \endlastfoot
\texttt{exp} & \texttt{x -{}-{} e\^{}x} & exponential (elementwise for matrices) & \texttt{1 exp}\\
\texttt{ln} & \texttt{x -{}-{} ln(x)} & natural logarithm; a negative real promotes to complex & \texttt{10 ln}\\
\texttt{log} & \texttt{x -{}-{} log10(x)} & base-{}10 logarithm; a negative real promotes to complex & \texttt{100 log}\\
\end{longtable}

\subsection*{Trigonometry}
\begin{longtable}{@{}>{\raggedright\arraybackslash}p{0.12\textwidth}>{\raggedright\arraybackslash}p{0.19\textwidth}>{\raggedright\arraybackslash}p{0.36\textwidth}>{\raggedright\arraybackslash}p{0.23\textwidth}@{}}
\toprule \textbf{Command} & \textbf{Stack} & \textbf{Description} & \textbf{Example}\\ \midrule \endfirsthead
\toprule \textbf{Command} & \textbf{Stack} & \textbf{Description} & \textbf{Example}\\ \midrule \endhead
\bottomrule \endlastfoot
\texttt{sin} & \texttt{z -{}-{} sin(z)} & sine of a real or complex argument (radians); elementwise for matrices & \texttt{0.5 sin}\\
\texttt{cos} & \texttt{z -{}-{} cos(z)} & cosine of real or complex argument & \texttt{1 cos}\\
\texttt{tan} & \texttt{z -{}-{} tan(z)} & tangent of real or complex argument & \texttt{0.3 tan}\\
\texttt{asin} & \texttt{z -{}-{} asin(z)} & inverse sine (principal value) & \texttt{0.5 asin}\\
\texttt{acos} & \texttt{z -{}-{} acos(z)} & inverse cosine (principal value) & \texttt{0.5 acos}\\
\texttt{atan} & \texttt{z -{}-{} atan(z)} & inverse tangent (principal value) & \texttt{1 atan}\\
\texttt{sinh} & \texttt{z -{}-{} sinh(z)} & hyperbolic sine & \texttt{1 sinh}\\
\texttt{cosh} & \texttt{z -{}-{} cosh(z)} & hyperbolic cosine & \texttt{1 cosh}\\
\texttt{tanh} & \texttt{z -{}-{} tanh(z)} & hyperbolic tangent & \texttt{1 tanh}\\
\texttt{asinh} & \texttt{z -{}-{} asinh(z)} & inverse hyperbolic sine & \texttt{1 asinh}\\
\texttt{acosh} & \texttt{z -{}-{} acosh(z)} & inverse hyperbolic cosine & \texttt{2 acosh}\\
\texttt{atanh} & \texttt{z -{}-{} atanh(z)} & inverse hyperbolic tangent & \texttt{0.5 atanh}\\
\end{longtable}

\subsection*{Complex}
\begin{longtable}{@{}>{\raggedright\arraybackslash}p{0.12\textwidth}>{\raggedright\arraybackslash}p{0.19\textwidth}>{\raggedright\arraybackslash}p{0.36\textwidth}>{\raggedright\arraybackslash}p{0.23\textwidth}@{}}
\toprule \textbf{Command} & \textbf{Stack} & \textbf{Description} & \textbf{Example}\\ \midrule \endfirsthead
\toprule \textbf{Command} & \textbf{Stack} & \textbf{Description} & \textbf{Example}\\ \midrule \endhead
\bottomrule \endlastfoot
\texttt{re} & \texttt{z -{}-{} Re(z)} & real part (a real is unchanged) & \texttt{(1,2) re}\\
\texttt{im} & \texttt{z -{}-{} Im(z)} & imaginary part (0 for a real) & \texttt{(1,2) im}\\
\texttt{arg} & \texttt{z -{}-{} arg(z)} & argument (phase) in radians, in (-{}pi, pi] & \texttt{(0,1) arg}\\
\texttt{conj} & \texttt{z -{}-{} conj(z)} & complex conjugate (elementwise for matrices; identity for reals) & \texttt{(1,2) conj}\\
\texttt{re2c} & \texttt{x -{}-{} (x,0)} & promote a real to complex & \texttt{5 re2c}\\
\texttt{split\_c} & \texttt{z -{}-{} Re(z) Im(z)} & split a complex number into real and imaginary parts & \texttt{(1,2) split\_c}\\
\texttt{j2r} & \texttt{re im -{}-{} (re,im)} & join two reals into a complex number (imaginary part on top) & \texttt{1 2 j2r}\\
\end{longtable}

\subsection*{Special functions}
\begin{longtable}{@{}>{\raggedright\arraybackslash}p{0.12\textwidth}>{\raggedright\arraybackslash}p{0.19\textwidth}>{\raggedright\arraybackslash}p{0.36\textwidth}>{\raggedright\arraybackslash}p{0.23\textwidth}@{}}
\toprule \textbf{Command} & \textbf{Stack} & \textbf{Description} & \textbf{Example}\\ \midrule \endfirsthead
\toprule \textbf{Command} & \textbf{Stack} & \textbf{Description} & \textbf{Example}\\ \midrule \endhead
\bottomrule \endlastfoot
\texttt{npdf} & \texttt{x -{}-{} phi(x)} & standard normal density & \texttt{0 npdf}\\
\texttt{ncdf} & \texttt{x -{}-{} Phi(x)} & standard normal cumulative distribution & \texttt{1.96 ncdf}\\
\texttt{nquant} & \texttt{p -{}-{} x} & standard normal quantile (inverse of ncdf), 0 \textless{} p \textless{} 1 & \texttt{0.975 nquant}\\
\texttt{gamma} & \texttt{x -{}-{} Gamma(x)} & gamma function & \texttt{5 gamma}\\
\texttt{ln\_gamma} & \texttt{x -{}-{} ln Γ(x)} & log gamma function & \texttt{10 ln\_gamma}\\
\texttt{beta} & \texttt{a b -{}-{} B(a,b)} & beta function & \texttt{2 3 beta}\\
\texttt{ln\_beta} & \texttt{a b -{}-{} ln B(a,b)} & log beta function & \texttt{2 3 ln\_beta}\\
\end{longtable}

\subsection*{Constants}
\begin{longtable}{@{}>{\raggedright\arraybackslash}p{0.12\textwidth}>{\raggedright\arraybackslash}p{0.19\textwidth}>{\raggedright\arraybackslash}p{0.36\textwidth}>{\raggedright\arraybackslash}p{0.23\textwidth}@{}}
\toprule \textbf{Command} & \textbf{Stack} & \textbf{Description} & \textbf{Example}\\ \midrule \endfirsthead
\toprule \textbf{Command} & \textbf{Stack} & \textbf{Description} & \textbf{Example}\\ \midrule \endhead
\bottomrule \endlastfoot
\texttt{pi} & \texttt{-{}-{} π} & push π & \texttt{pi}\\
\texttt{e} & \texttt{-{}-{} e} & push Euler\textquotesingle{}s number & \texttt{e}\\
\texttt{gravity} & \texttt{-{}-{} g} & standard gravity, 9.80665 m/s\^{}2 & \texttt{gravity}\\
\texttt{inf} & \texttt{-{}-{} +∞} & positive infinity & \texttt{inf}\\
\texttt{nan} & \texttt{-{}-{} NaN} & not-{}a-{}number & \texttt{nan}\\
\end{longtable}

\subsection*{Logic}
\begin{longtable}{@{}>{\raggedright\arraybackslash}p{0.12\textwidth}>{\raggedright\arraybackslash}p{0.19\textwidth}>{\raggedright\arraybackslash}p{0.36\textwidth}>{\raggedright\arraybackslash}p{0.23\textwidth}@{}}
\toprule \textbf{Command} & \textbf{Stack} & \textbf{Description} & \textbf{Example}\\ \midrule \endfirsthead
\toprule \textbf{Command} & \textbf{Stack} & \textbf{Description} & \textbf{Example}\\ \midrule \endhead
\bottomrule \endlastfoot
\texttt{eq} & \texttt{a b -{}-{} flag} & 1 if a equals b, else 0; numbers, strings, or elementwise for matrices & \texttt{2 2 eq}\\
\texttt{neq} & \texttt{a b -{}-{} flag} & 1 if a differs from b; numbers, strings, or elementwise for matrices & \texttt{2 3 neq}\\
\texttt{lt} & \texttt{a b -{}-{} flag} & 1 if a \textless{} b (complex numbers compare by magnitude); elementwise for matrices & \texttt{2 3 lt}\\
\texttt{leq} & \texttt{a b -{}-{} flag} & 1 if a \textless{}= b; elementwise for matrices & \texttt{2 3 leq}\\
\texttt{gt} & \texttt{a b -{}-{} flag} & 1 if a \textgreater{} b; elementwise for matrices & \texttt{3 2 gt}\\
\texttt{geq} & \texttt{a b -{}-{} flag} & 1 if a \textgreater{}= b; elementwise for matrices & \texttt{3 2 geq}\\
\texttt{and} & \texttt{a b -{}-{} flag} & logical and (0 is false, anything else true) & \texttt{1 0 and}\\
\texttt{or} & \texttt{a b -{}-{} flag} & logical or & \texttt{1 0 or}\\
\texttt{not} & \texttt{a -{}-{} flag} & logical not & \texttt{0 not}\\
\end{longtable}

\subsection*{Stack}
\begin{longtable}{@{}>{\raggedright\arraybackslash}p{0.12\textwidth}>{\raggedright\arraybackslash}p{0.19\textwidth}>{\raggedright\arraybackslash}p{0.36\textwidth}>{\raggedright\arraybackslash}p{0.23\textwidth}@{}}
\toprule \textbf{Command} & \textbf{Stack} & \textbf{Description} & \textbf{Example}\\ \midrule \endfirsthead
\toprule \textbf{Command} & \textbf{Stack} & \textbf{Description} & \textbf{Example}\\ \midrule \endhead
\bottomrule \endlastfoot
\texttt{dup} & \texttt{x -{}-{} x x} & duplicate top of stack & \texttt{5 dup}\\
\texttt{drop} & \texttt{x -{}-{}} & drop the top stack element & \texttt{1 2 drop}\\
\texttt{swap} & \texttt{a b -{}-{} b a} & swap top two elements & \texttt{1 2 swap}\\
\texttt{over} & \texttt{a b -{}-{} a b a} & copy second item to top & \texttt{1 2 over}\\
\texttt{nip} & \texttt{a b -{}-{} b} & drop second-{}from-{}top element (keep top) & \texttt{1 2 nip}\\
\texttt{tuck} & \texttt{a b -{}-{} b a b} & duplicate second item and tuck below top & \texttt{1 2 tuck}\\
\texttt{roll} & \texttt{... n -{}-{} ...} & move the item at depth n to the top (1 roll = swap, 2 roll = rot) & \texttt{1 2 3 2 roll}\\
\texttt{pick} & \texttt{... n -{}-{} ... v} & copy the item at depth n to the top (0 pick = dup) & \texttt{1 2 3 2 pick}\\
\texttt{depth} & \texttt{-{}-{} n} & number of items on the stack & \texttt{depth}\\
\texttt{clst} & \texttt{... -{}-{}} & clear the stack & \texttt{clst}\\
\texttt{savestack} & \texttt{-{}-{}} & save the whole stack to stack\_backup.dat in the config directory & \texttt{savestack}\\
\texttt{loadstack} & \texttt{-{}-{} ...} & push the stack saved by savestack & \texttt{loadstack}\\
\end{longtable}

\subsection*{Strings}
\begin{longtable}{@{}>{\raggedright\arraybackslash}p{0.12\textwidth}>{\raggedright\arraybackslash}p{0.19\textwidth}>{\raggedright\arraybackslash}p{0.36\textwidth}>{\raggedright\arraybackslash}p{0.23\textwidth}@{}}
\toprule \textbf{Command} & \textbf{Stack} & \textbf{Description} & \textbf{Example}\\ \midrule \endfirsthead
\toprule \textbf{Command} & \textbf{Stack} & \textbf{Description} & \textbf{Example}\\ \midrule \endhead
\bottomrule \endlastfoot
\texttt{scon} & \texttt{s1 s2 -{}-{} s1s2} & concatenate two strings & \texttt{\textquotedbl{}foo\textquotedbl{} \textquotedbl{}bar\textquotedbl{} scon}\\
\texttt{substr} & \texttt{s start end -{}-{} sub} & substring from index start (0-{}based) up to end (exclusive) & \texttt{\textquotedbl{}abcdef\textquotedbl{} 1 3 substr}\\
\texttt{s2l} & \texttt{s -{}-{} s\_lower} & convert string to lowercase & \texttt{\textquotedbl{}Hello\textquotedbl{} s2l}\\
\texttt{s2u} & \texttt{s -{}-{} s\_upper} & convert string to uppercase & \texttt{\textquotedbl{}Hello\textquotedbl{} s2u}\\
\texttt{slen} & \texttt{s -{}-{} n} & length of a string in bytes & \texttt{\textquotedbl{}hello\textquotedbl{} slen}\\
\texttt{srev} & \texttt{s -{}-{} s\_rev} & reverse string & \texttt{\textquotedbl{}abc\textquotedbl{} srev}\\
\texttt{int2str} & \texttt{x -{}-{} s} & integer part of x as a string & \texttt{123 int2str}\\
\end{longtable}

\subsection*{Matrices}
\begin{longtable}{@{}>{\raggedright\arraybackslash}p{0.12\textwidth}>{\raggedright\arraybackslash}p{0.19\textwidth}>{\raggedright\arraybackslash}p{0.36\textwidth}>{\raggedright\arraybackslash}p{0.23\textwidth}@{}}
\toprule \textbf{Command} & \textbf{Stack} & \textbf{Description} & \textbf{Example}\\ \midrule \endfirsthead
\toprule \textbf{Command} & \textbf{Stack} & \textbf{Description} & \textbf{Example}\\ \midrule \endhead
\bottomrule \endlastfoot
\texttt{tran} & \texttt{A -{}-{} A\^{}T} & transpose & \texttt{[2 3 \$ 1 2 3 4 5 6] tran}\\
\texttt{\textquotesingle{}} & \texttt{A -{}-{} A\textquotesingle{}} & transpose (alias of tran) & \texttt{[2 3 \$ 1 2 3 4 5 6] \textquotesingle{}}\\
\texttt{minv} & \texttt{A -{}-{} A\^{}-{}1} & inverse of a square real or complex matrix (error if singular) & \texttt{[2 2 \$ 4 7 2 6] minv}\\
\texttt{pinv} & \texttt{A -{}-{} A\^{}+} & Moore-{}Penrose pseudoinverse of a square real matrix (SVD) & \texttt{[2 2 \$ 1 2 2 4] pinv}\\
\texttt{det} & \texttt{A -{}-{} det(A)} & determinant of a square real or complex matrix & \texttt{[2 2 \$ 4 7 2 6] det}\\
\texttt{eig} & \texttt{A -{}-{} V L} & eigenvectors (columns of V) and eigenvalues (diagonal of L) of a square real matrix; both complex & \texttt{[2 2 \$ 2 1 1 2] eig}\\
\texttt{chol} & \texttt{A -{}-{} L} & Cholesky factor of a symmetric positive definite real matrix, A = L L\^{}T & \texttt{[2 2 \$ 4 2 2 3] chol}\\
\texttt{svd} & \texttt{A -{}-{} U S V} & singular value decomposition of a real matrix, A = U S V\^{}T & \texttt{[2 2 \$ 4 7 2 6] svd}\\
\texttt{kron} & \texttt{A B -{}-{} A(x)B} & Kronecker product & \texttt{[2 2 \$ 1 2 3 4] [1 2 \$ 5 6] kron}\\
\texttt{dim} & \texttt{A -{}-{} A rows cols} & push the dimensions (the matrix stays) & \texttt{[2 3 \$ 1 2 3 4 5 6] dim}\\
\texttt{reshape} & \texttt{A rows cols -{}-{} B} & reshape in row-{}major order; rows*cols must match & \texttt{[2 3 \$ 1 2 3 4 5 6] 3 2 reshape}\\
\texttt{get\_aij} & \texttt{A i j -{}-{} A a\_ij} & push element (i,j), 0-{}based (the matrix stays) & \texttt{[2 2 \$ 1 2 3 4] 0 1 get\_aij}\\
\texttt{set\_aij} & \texttt{A i j x -{}-{} A} & set element (i,j), 0-{}based, in place; a complex x promotes a real matrix & \texttt{[2 2 \$ 1 2 3 4] 0 1 9 set\_aij}\\
\texttt{split\_mat} & \texttt{A -{}-{} a11 a12 ...} & push all elements as scalars, row by row & \texttt{[2 2 \$ 1 2 3 4] split\_mat}\\
\texttt{diag} & \texttt{A -{}-{} v} & main diagonal as a row vector & \texttt{[2 2 \$ 1 2 3 4] diag}\\
\texttt{to\_diag} & \texttt{v -{}-{} D} & diagonal matrix from a row or column vector & \texttt{[1 3 \$ 1 2 3] to\_diag}\\
\texttt{eye} & \texttt{n -{}-{} I} & n x n identity matrix & \texttt{3 eye}\\
\texttt{ones} & \texttt{rows cols -{}-{} A} & matrix filled with ones & \texttt{2 3 ones}\\
\texttt{zeroes} & \texttt{rows cols -{}-{} A} & matrix filled with zeros & \texttt{2 3 zeroes}\\
\texttt{rand} & \texttt{rows cols -{}-{} A} & matrix of uniform(0,1) randoms & \texttt{2 2 rand}\\
\texttt{randn} & \texttt{rows cols -{}-{} A} & matrix of standard normal randoms & \texttt{2 2 randn}\\
\texttt{rrange} & \texttt{start step end -{}-{} v} & row vector start, start+step, ... up to end & \texttt{0 0.25 1 rrange}\\
\texttt{join\_v} & \texttt{A B -{}-{} [A; B]} & stack two matrices vertically (same number of columns) & \texttt{[1 2 \$ 1 2] [1 2 \$ 3 4] join\_v}\\
\texttt{join\_h} & \texttt{A B -{}-{} [A B]} & join two matrices side by side (same number of rows) & \texttt{[2 1 \$ 1 2] [2 1 \$ 3 4] join\_h}\\
\texttt{cumsum\_r} & \texttt{A -{}-{} B} & cumulative sum along each row & \texttt{[2 2 \$ 1 2 3 4] cumsum\_r}\\
\texttt{cumsum\_c} & \texttt{A -{}-{} B} & cumulative sum down each column & \texttt{[2 2 \$ 1 2 3 4] cumsum\_c}\\
\texttt{cmean} & \texttt{A -{}-{} v} & mean of each column (row vector) & \texttt{[2 2 \$ 1 2 3 4] cmean}\\
\texttt{rmean} & \texttt{A -{}-{} v} & mean of each row (column vector) & \texttt{[2 2 \$ 1 2 3 4] rmean}\\
\texttt{csum} & \texttt{A -{}-{} v} & sum of each column (row vector) & \texttt{[2 2 \$ 1 2 3 4] csum}\\
\texttt{rsum} & \texttt{A -{}-{} v} & sum of each row (column vector) & \texttt{[2 2 \$ 1 2 3 4] rsum}\\
\texttt{cvar} & \texttt{A -{}-{} v} & sample variance of each column (row vector) & \texttt{[2 2 \$ 1 2 3 4] cvar}\\
\texttt{rvar} & \texttt{A -{}-{} v} & sample variance of each row (column vector) & \texttt{[2 2 \$ 1 2 3 4] rvar}\\
\texttt{cmin} & \texttt{A -{}-{} v} & minimum of each column & \texttt{[2 2 \$ 1 2 3 4] cmin}\\
\texttt{rmin} & \texttt{A -{}-{} v} & minimum of each row & \texttt{[2 2 \$ 1 2 3 4] rmin}\\
\texttt{cmax} & \texttt{A -{}-{} v} & maximum of each column & \texttt{[2 2 \$ 1 2 3 4] cmax}\\
\texttt{rmax} & \texttt{A -{}-{} v} & maximum of each row & \texttt{[2 2 \$ 1 2 3 4] rmax}\\
\end{longtable}

\subsection*{Polynomials}
\begin{longtable}{@{}>{\raggedright\arraybackslash}p{0.12\textwidth}>{\raggedright\arraybackslash}p{0.19\textwidth}>{\raggedright\arraybackslash}p{0.36\textwidth}>{\raggedright\arraybackslash}p{0.23\textwidth}@{}}
\toprule \textbf{Command} & \textbf{Stack} & \textbf{Description} & \textbf{Example}\\ \midrule \endfirsthead
\toprule \textbf{Command} & \textbf{Stack} & \textbf{Description} & \textbf{Example}\\ \midrule \endhead
\bottomrule \endlastfoot
\texttt{roots} & \texttt{coeffs -{}-{} roots} & roots of a real polynomial; coefficients highest degree first, result is a complex row vector & \texttt{[1 3 \$ 1 -{}3 2] roots}\\
\texttt{pval} & \texttt{coeffs x -{}-{} y} & evaluate a polynomial (coefficients highest degree first, real or complex) at x & \texttt{[1 3 \$ 1 0 -{}1] 2 pval}\\
\end{longtable}

\subsection*{Numerics}
\begin{longtable}{@{}>{\raggedright\arraybackslash}p{0.12\textwidth}>{\raggedright\arraybackslash}p{0.19\textwidth}>{\raggedright\arraybackslash}p{0.36\textwidth}>{\raggedright\arraybackslash}p{0.23\textwidth}@{}}
\toprule \textbf{Command} & \textbf{Stack} & \textbf{Description} & \textbf{Example}\\ \midrule \endfirsthead
\toprule \textbf{Command} & \textbf{Stack} & \textbf{Description} & \textbf{Example}\\ \midrule \endhead
\bottomrule \endlastfoot
\texttt{exec} & \texttt{xt -{}-{}} & run an execution token (@name pushes one) & \texttt{3 @sq exec}\\
\texttt{integrate} & \texttt{a b xt -{}-{} I} & integral of f over [a,b] (tanh-{}sinh with adaptive bisection; handles endpoint singularities) & \texttt{0 pi @sin integrate}\\
\texttt{romberg} & \texttt{a b xt -{}-{} I} & integral of f over [a,b] by Romberg extrapolation & \texttt{0 1 @exp romberg}\\
\texttt{fzero} & \texttt{a b xt -{}-{} x} & root of f in the bracket [a,b] where f changes sign (Brent) & \texttt{2 4 @sin fzero}\\
\texttt{solve} & \texttt{guess xt -{}-{} x} & root of f near one guess (bracket search, then Brent) & \texttt{3 @sin solve}\\
\texttt{minimize} & \texttt{a b xt -{}-{} x} & argument of the minimum of f in [a,b] (Brent) & \texttt{0 5 @par minimize}\\
\texttt{deriv} & \texttt{x xt -{}-{} f\textquotesingle{}(x)} & derivative of f at x (Ridders extrapolation) & \texttt{2 @exp deriv}\\
\texttt{fsum} & \texttt{lo hi xt -{}-{} S} & sum of f(i) for integer i from lo to hi & \texttt{1 1000 @invsq fsum}\\
\texttt{fprod} & \texttt{lo hi xt -{}-{} P} & product of f(i) for integer i from lo to hi & \texttt{1 10 @ident fprod}\\
\texttt{tol} & \texttt{x -{}-{}} & set the convergence tolerance (default 1e-{}10) & \texttt{1e-{}6 tol}\\
\texttt{nevals} & \texttt{-{}-{} n} & function evaluations used by the last numerical call & \texttt{nevals}\\
\texttt{interr} & \texttt{-{}-{} e} & error estimate of the last integrate & \texttt{interr}\\
\end{longtable}

\subsection*{Registers}
\begin{longtable}{@{}>{\raggedright\arraybackslash}p{0.12\textwidth}>{\raggedright\arraybackslash}p{0.19\textwidth}>{\raggedright\arraybackslash}p{0.36\textwidth}>{\raggedright\arraybackslash}p{0.23\textwidth}@{}}
\toprule \textbf{Command} & \textbf{Stack} & \textbf{Description} & \textbf{Example}\\ \midrule \endfirsthead
\toprule \textbf{Command} & \textbf{Stack} & \textbf{Description} & \textbf{Example}\\ \midrule \endhead
\bottomrule \endlastfoot
\texttt{sto} & \texttt{x n -{}-{}} & store x in register n (0-{}63); any type & \texttt{42 0 sto}\\
\texttt{rcl} & \texttt{n -{}-{} x} & push a copy of register n & \texttt{0 rcl}\\
\texttt{ffr} & \texttt{-{}-{} n} & index of the first free register & \texttt{ffr}\\
\texttt{pr} & \texttt{-{}-{}} & show which registers are occupied & \texttt{pr}\\
\texttt{saveregs} & \texttt{-{}-{}} & save all registers to registers.txt in the config directory & \texttt{saveregs}\\
\texttt{loadregs} & \texttt{-{}-{}} & load the registers saved by saveregs & \texttt{loadregs}\\
\texttt{clregs} & \texttt{-{}-{}} & clear all registers & \texttt{clregs}\\
\end{longtable}

\subsection*{Dates}
\begin{longtable}{@{}>{\raggedright\arraybackslash}p{0.12\textwidth}>{\raggedright\arraybackslash}p{0.19\textwidth}>{\raggedright\arraybackslash}p{0.36\textwidth}>{\raggedright\arraybackslash}p{0.23\textwidth}@{}}
\toprule \textbf{Command} & \textbf{Stack} & \textbf{Description} & \textbf{Example}\\ \midrule \endfirsthead
\toprule \textbf{Command} & \textbf{Stack} & \textbf{Description} & \textbf{Example}\\ \midrule \endhead
\bottomrule \endlastfoot
\texttt{today} & \texttt{-{}-{} date} & today\textquotesingle{}s date as a string d.m.yyyy & \texttt{today}\\
\texttt{now} & \texttt{-{}-{} \textquotedbl{}HH:MM\textquotedbl{}} & push the current local time as a string, same form as sunrise/sunset & \texttt{now}\\
\texttt{tz\_offset} & \texttt{date zone -{}-{} hours} & UTC offset in hours of an IANA zone on the given date, DST included & \texttt{\textquotedbl{}1.12.2026\textquotedbl{} \textquotedbl{}America/New\_York\textquotedbl{} tz\_offset}\\
\texttt{ddays} & \texttt{d1 d2 -{}-{} n} & days from date d1 to date d2 (d2 -{} d1) & \texttt{\textquotedbl{}1.1.2025\textquotedbl{} \textquotedbl{}31.12.2025\textquotedbl{} ddays}\\
\texttt{dateplus} & \texttt{date n -{}-{} date2} & the date n days later (n may be negative) & \texttt{\textquotedbl{}1.1.2025\textquotedbl{} 30 dateplus}\\
\texttt{dow} & \texttt{date -{}-{} name} & weekday name of a date & \texttt{\textquotedbl{}1.1.2025\textquotedbl{} dow}\\
\texttt{edmy} & \texttt{date -{}-{} y m d} & split a date into year, month, day (day on top) & \texttt{\textquotedbl{}1.1.2025\textquotedbl{} edmy}\\
\texttt{num2date} & \texttt{d m y -{}-{} date} & build a date string from day, month, year & \texttt{1 1 2025 num2date}\\
\texttt{days2eoy} & \texttt{date -{}-{} n} & days from the date to the end of its year & \texttt{\textquotedbl{}1.1.2025\textquotedbl{} days2eoy}\\
\texttt{sunrise} & \texttt{date lat lon zone -{}-{} time} & local sunrise HH:MM; zone is an IANA name or a UTC offset in hours & \texttt{\textquotedbl{}21.6.2025\textquotedbl{} 38.9 -{}77.0 \textquotedbl{}America/New\_York\textquotedbl{} sunrise}\\
\texttt{sunset} & \texttt{date lat lon zone -{}-{} time} & local sunset HH:MM & \texttt{\textquotedbl{}21.6.2025\textquotedbl{} 38.9 -{}77.0 -{}4 sunset}\\
\texttt{dawn} & \texttt{date lat lon zone -{}-{} time} & civil dawn HH:MM (sun 6 degrees below the horizon) & \texttt{today 38.9 -{}77.0 \textquotedbl{}America/New\_York\textquotedbl{} dawn}\\
\texttt{dusk} & \texttt{date lat lon zone -{}-{} time} & civil dusk HH:MM & \texttt{today 38.9 -{}77.0 \textquotedbl{}America/New\_York\textquotedbl{} dusk}\\
\end{longtable}

\subsection*{Display}
\begin{longtable}{@{}>{\raggedright\arraybackslash}p{0.12\textwidth}>{\raggedright\arraybackslash}p{0.19\textwidth}>{\raggedright\arraybackslash}p{0.36\textwidth}>{\raggedright\arraybackslash}p{0.23\textwidth}@{}}
\toprule \textbf{Command} & \textbf{Stack} & \textbf{Description} & \textbf{Example}\\ \midrule \endfirsthead
\toprule \textbf{Command} & \textbf{Stack} & \textbf{Description} & \textbf{Example}\\ \midrule \endhead
\bottomrule \endlastfoot
\texttt{print} & \texttt{x -{}-{} x} & print the top of the stack (a matrix in full) & \texttt{42 print}\\
\texttt{pm} & \texttt{A -{}-{} A} & print the matrix on top of the stack in full & \texttt{[2 2 \$ 1 2 3 4] pm}\\
\texttt{ps} & \texttt{-{}-{}} & print the whole stack & \texttt{ps}\\
\texttt{setprec} & \texttt{n -{}-{}} & digits after the decimal point (fixed) or significant digits (scientific), 0-{}20 & \texttt{10 setprec}\\
\texttt{sfs} & \texttt{-{}-{}} & toggle between fixed-{}point and scientific display & \texttt{sfs}\\
\end{longtable}

\subsection*{Words}
\begin{longtable}{@{}>{\raggedright\arraybackslash}p{0.12\textwidth}>{\raggedright\arraybackslash}p{0.19\textwidth}>{\raggedright\arraybackslash}p{0.36\textwidth}>{\raggedright\arraybackslash}p{0.23\textwidth}@{}}
\toprule \textbf{Command} & \textbf{Stack} & \textbf{Description} & \textbf{Example}\\ \midrule \endfirsthead
\toprule \textbf{Command} & \textbf{Stack} & \textbf{Description} & \textbf{Example}\\ \midrule \endhead
\bottomrule \endlastfoot
\texttt{listwords} & \texttt{-{}-{}} & list the user-{}defined words & \texttt{listwords}\\
\texttt{listmacros} & \texttt{-{}-{}} & list the macros loaded from predefined\_macros.txt & \texttt{listmacros}\\
\texttt{loadwords} & \texttt{-{}-{}} & load the words saved in user\_words.txt (also done at startup) & \texttt{loadwords}\\
\texttt{savewords} & \texttt{-{}-{}} & save all user-{}defined words to user\_words.txt in the config directory & \texttt{savewords}\\
\texttt{delword} & \texttt{\textquotedbl{}name\textquotedbl{} -{}-{}} & delete a user-{}defined word & \texttt{\textquotedbl{}sq\textquotedbl{} delword}\\
\texttt{clrwords} & \texttt{-{}-{}} & delete all user-{}defined words & \texttt{clrwords}\\
\end{longtable}

\subsection*{Programs}
\begin{longtable}{@{}>{\raggedright\arraybackslash}p{0.12\textwidth}>{\raggedright\arraybackslash}p{0.19\textwidth}>{\raggedright\arraybackslash}p{0.36\textwidth}>{\raggedright\arraybackslash}p{0.23\textwidth}@{}}
\toprule \textbf{Command} & \textbf{Stack} & \textbf{Description} & \textbf{Example}\\ \midrule \endfirsthead
\toprule \textbf{Command} & \textbf{Stack} & \textbf{Description} & \textbf{Example}\\ \midrule \endhead
\bottomrule \endlastfoot
\texttt{eval} & \texttt{s -{}-{}} & evaluate the string s as a line of input & \texttt{\textquotedbl{}1 2 +\textquotedbl{} eval}\\
\texttt{batch} & \texttt{s -{}-{}} & run the lines of file s as input & \texttt{\textquotedbl{}calc.txt\textquotedbl{} batch}\\
\texttt{run} & \texttt{s -{}-{}} & load and run the HP-{}41 style program in file s (LBL/GOTO/GOSUB/RTN/END) & \texttt{\textquotedbl{}prog1.prg\textquotedbl{} run}\\
\end{longtable}

\subsection*{Session}
\begin{longtable}{@{}>{\raggedright\arraybackslash}p{0.12\textwidth}>{\raggedright\arraybackslash}p{0.19\textwidth}>{\raggedright\arraybackslash}p{0.36\textwidth}>{\raggedright\arraybackslash}p{0.23\textwidth}@{}}
\toprule \textbf{Command} & \textbf{Stack} & \textbf{Description} & \textbf{Example}\\ \midrule \endfirsthead
\toprule \textbf{Command} & \textbf{Stack} & \textbf{Description} & \textbf{Example}\\ \midrule \endhead
\bottomrule \endlastfoot
\texttt{help} & \texttt{-{}-{}} & print the quick-{}start help screen & \texttt{help}\\
\texttt{usage} & \texttt{\textquotedbl{}name\textquotedbl{} -{}-{}} & show the stack picture, description and example of a command, or a word\textquotesingle{}s definition & \texttt{\textquotedbl{}sin\textquotedbl{} usage}\\
\texttt{listfcns} & \texttt{-{}-{}} & list all commands by category & \texttt{listfcns}\\
\texttt{clrhist} & \texttt{-{}-{}} & clear the readline history & \texttt{clrhist}\\
\texttt{fuck} & \texttt{-{}-{}} & an Easter egg & \texttt{fuck}\\
\end{longtable}


\normalsize

\clearpage
\section{Shipped macros}
\label{app:macros}

The following words are loaded from \code{predefined\_macros.txt}. The file is replaced at each \code{make install}; add your own words to \code{user\_words.txt} instead.

\begin{longtable}{@{}>{\raggedright\arraybackslash}p{0.13\textwidth}>{\raggedright\arraybackslash}p{0.38\textwidth}>{\raggedright\arraybackslash}p{0.41\textwidth}@{}}
\toprule \textbf{Macro} & \textbf{Definition} & \textbf{Meaning} \\ \midrule \endfirsthead
\toprule \textbf{Macro} & \textbf{Definition} & \textbf{Meaning} \\ \midrule \endhead
\bottomrule \endlastfoot
\code{sq}      & \code{dup *}            & $x^2$ \\
\code{qu}      & \code{sq sq}            & $x^4$ \\
\code{cube}    & \code{3 \textasciicircum} & $x^3$ \\
\code{cubrt}   & \code{3 inv \textasciicircum} & cube root (complex for negative $x$) \\
\code{pctchg}  & \code{over - swap / 100 *} & percent change from $y$ to $x$, consuming both \\
\code{pctot}   & \code{/ 100 *}          & $x$ as a percentage of $y$, consuming both \\
\code{ln2}     & \code{ln 2 ln /}        & base-2 logarithm \\
\code{norm2}   & \code{svd drop swap drop diag rmax split\_mat} & spectral norm (largest singular value) \\
\code{normfro} & \code{dim * 1 reshape 2 .\textasciicircum{} csum split\_mat sqrt} & Frobenius norm \\
\code{cond}    & \code{svd drop swap drop diag dup rmax swap rmin ./ split\_mat} & 2-norm condition number \\
\code{isnan}, \code{isinf} & \code{nan eq}, \code{inf eq} & flags \\
\code{isreal}  & \code{im 0 eq}          & 1 if the imaginary part is zero \\
\code{selreal} & \code{dup isreal swap re .*} & keep the real entries of a vector, zero the rest \\
\code{npv}     & \code{1 + inv pval}     & net present value, \code{cashflows rate npv}; the cash flows are a row vector with the \emph{latest} period first, so \code{[1 3 \$ 121 0 -100] 0.1 npv} is 0 \\
\code{irr}     & \code{roots selreal rmax split\_mat inv 1 - 100 *} & internal rate of return in percent, same cash-flow convention: \code{[1 3 \$ 121 0 -100] irr} is 10 \\
\code{cumprod\_c}, \code{cumprod\_r} & \code{ln cumsum\_c exp} & cumulative products \\
\code{sto\_add}, \code{sto\_sub}, \code{sto\_mul}, \code{sto\_div} & \code{dup rcl 2 roll + swap sto} etc. & storage arithmetic: \code{x n sto\_add} adds $x$ to register $n$ \\
\code{XprimeX} & \code{dup ' swap *}     & $X^T X$ \\
\code{f2c}, \code{c2f} & \code{32 - 9 / 5 *}, \code{9 * 5 / 32 +} & Fahrenheit and Celsius \\
\code{dec\_hr} & \code{dup 0 2 substr eval swap 3 5 substr eval 60 / +} & \code{"HH:MM"} to decimal hours \\
\code{alexandria}, \code{duluth}, \code{durham}, \code{kranj} & \code{lat lon "zone"} & locations for the solar functions \\
\end{longtable}

\clearpage
\section{Changes from version 1.x}
\label{app:changes}

Users of MM-15 1.x will notice these differences; \code{CHANGES.md} in the source tree has the complete list.
\begin{itemize}
  \item Errors are atomic and roll the stack back; batch mode exits 1 on error.
  \item One promotion rule for \cmd{sqrt}, \cmd{ln}, \cmd{log}, \cmd{asin}, \cmd{acos}, \cmd{acosh}, \cmd{atanh}, \cmd{\textasciicircum}; division by zero is an error.
  \item \cmd{sto} takes \stk{x n} (HP order); \cmd{substr} takes \stk{s start end}; \cmd{set\_aij} takes \stk{A i j x} like \cmd{get\_aij}; \cmd{roll} honours its argument; \cmd{pick} and \cmd{depth} are new.
  \item \cmd{roots} takes coefficients highest degree first, consistent with \cmd{pval}.
  \item \cmd{rrange} takes \stk{start step end}; \cmd{days2eoy} takes a date; \cmd{eq} and \cmd{neq} compare strings; real and complex matrices mix in arithmetic.
  \item Numerical methods take \code{@name} tokens: \code{selword}, \code{set\_intg\_tol} and \code{set\_f0\_tol} are gone, \cmd{tol} replaces them; \cmd{romberg}, \cmd{solve}, \cmd{minimize}, \cmd{deriv}, \cmd{fsum}, \cmd{fprod}, \cmd{exec} are new.
  \item Words are compiled at definition time and have \code{if/else/then}, \code{begin/until}, \code{begin/while/repeat}, \code{do/loop/+loop}.
  \item Startup no longer makes network requests; the configuration directory is found at run time.
\end{itemize}

\end{document}
